\section{Execution Unit dan Urutan Pengambilan--Eksekusi}
\label{sec:eu}

\eng{Execution Unit} (EU) adalah bagian prosesor yang menerjemahkan
dan menjalankan instruksi. Jika BIU diibaratkan sebagai pintu dan
jalur pengiriman, EU adalah tempat pekerjaan utama dilakukan. Bagian
ini diringkas dari \cite{gunawan2019}, Bagian 2, hlm.~15--18.

\subsection{Tugas Execution Unit}

EU mempunyai tiga tugas utama:

\begin{enumerate}
    \item menerjemahkan perintah yang diterima dari BIU;
    \item mengeksekusi perintah tersebut;
    \item memberi tahu BIU dari mana instruksi atau data harus
          diambil.
\end{enumerate}

EU tidak berhubungan langsung dengan bus luar. Ketika instruksi
memerlukan operand dari memori, EU menyampaikan alamat yang
diperlukan kepada BIU, lalu menunggu data tersebut dikirim kembali.

\subsection{Komponen Execution Unit}

Komponen EU dan fungsinya dirangkum pada Tabel~\ref{tab:komponen-eu}.

\begin{table}[H]
    \centering
    \caption{Komponen utama Execution Unit dan fungsinya.}
    \label{tab:komponen-eu}
    \small
    \begin{tabular}{@{}p{4.2cm}p{9.0cm}@{}}
        \toprule
        \textbf{Komponen} & \textbf{Fungsi} \\
        \midrule
        ALU & Menjalankan operasi aritmetika dan logika, seperti
        penjumlahan, pengurangan, perkalian, pembagian, AND, OR, dan
        NOT \\
        \addlinespace
        Register \eng{flag} & Menyimpan penanda hasil operasi ALU,
        misalnya hasil nol, hasil negatif, atau terjadi \eng{carry} \\
        \addlinespace
        \eng{Operand} & Data yang dioperasikan; dapat berupa nilai,
        register, atau variabel \\
        \addlinespace
        \eng{Instruction decoder and control} & Menerjemahkan
        perintah dari BIU dan mengendalikan jalannya operasi di
        dalam EU \\
        \addlinespace
        Register umum & Menyimpan data yang sedang diolah; terdiri
        atas AX, BX, CX, dan DX \\
        \addlinespace
        Register penunjuk dan indeks & SP, BP, SI, dan DI; menunjuk
        lokasi di memori untuk \eng{stack}, data, atau operasi string \\
        \bottomrule
    \end{tabular}
\end{table}

Register umum 16 bit dapat dipisah menjadi dua bagian 8 bit. AX
terdiri atas AH dan AL, BX terdiri atas BH dan BL, dan seterusnya.
Akses per byte ini penting ketika program mengolah data 8 bit.
Pembahasan rinci setiap register disajikan pada bab tentang register
dan \eng{flag}.

\subsection{ALU, Flag, dan Operand}

ALU menjalankan operasi yang berhubungan dengan aritmetika dan
logika. Hasil operasi ditandai melalui \eng{flag}. Sebagai contoh,
\eng{flag} nol menandakan hasil operasi bernilai nol, sedangkan
\eng{flag} \eng{carry} menandakan hasil penjumlahan melebihi daya
tampung register yang dipakai.

\eng{Operand} adalah data yang dioperasikan. Pada instruksi
\lstasm|ADD AL, 5|, \lstasm|AL| dan \lstasm|5| adalah operand,
sedangkan \lstasm|ADD| adalah perintah yang menentukan operasi.
Operand dapat berupa nilai langsung, register, atau isi memori.

\subsection{Urutan Pengambilan dan Eksekusi}

Satu instruksi melewati tiga tahap utama seperti ditunjukkan pada
Gambar~\ref{fig:siklus-instruksi}.

\begin{figure}[H]
    \centering
    \begin{tikzpicture}[
        tahap/.style={
            draw=obeBlue,
            line width=0.9pt,
            fill=obeBlueLight,
            rounded corners=3pt,
            minimum width=2.9cm,
            minimum height=1.2cm,
            align=center,
            font=\small\bfseries
        },
        panah/.style={-{Stealth[length=2.5mm]}, line width=1pt, obeTeal},
        unit/.style={font=\scriptsize, text=obeGray}
    ]
        \node[tahap] (f) {Fetch\\{\normalfont\scriptsize ambil instruksi}};
        \node[tahap, right=0.8cm of f, fill=obeTealLight,
              draw=obeTeal] (d) {Decode\\{\normalfont\scriptsize terjemahkan}};
        \node[tahap, right=0.8cm of d, fill=obeAmberLight,
              draw=obeAmber] (e) {Execute\\{\normalfont\scriptsize jalankan}};

        \draw[panah] (f) -- (d);
        \draw[panah] (d) -- (e);

        \node[unit, below=0.12cm of f] {dikerjakan BIU};
        \node[unit, below=0.12cm of d] {dikerjakan EU};
        \node[unit, below=0.12cm of e] {dikerjakan EU};
    \end{tikzpicture}
    \caption{Tiga tahap utama pemrosesan instruksi pada 8086 beserta
    unit yang mengerjakannya.}
    \label{fig:siklus-instruksi}
\end{figure}

\begin{enumerate}
    \item \textbf{Fetch.} BIU menghitung alamat fisik dari
          \reg{CS}:\reg{IP}, mengambil byte instruksi melalui bus,
          lalu memasukkannya ke \eng{queue}. \reg{IP} bertambah
          sesuai jumlah byte yang diambil.

    \item \textbf{Decode.} EU mengambil byte instruksi dari
          \eng{queue} dan menerjemahkannya. EU menentukan operasi
          apa yang harus dilakukan serta operand apa saja yang
          terlibat.

    \item \textbf{Execute.} EU menjalankan operasi melalui ALU atau
          memindahkan data sesuai instruksi. Bila operand berada di
          memori, EU meminta BIU membacanya terlebih dahulu.
\end{enumerate}

Hasil eksekusi dapat ditulis ke register atau ke memori. Penulisan
ke register tidak memerlukan bus luar, sedangkan penulisan ke memori
memerlukan siklus bus tambahan melalui BIU. Perbedaan biaya inilah
yang akan diamati pada contoh dan praktikum bab ini.

\begin{infobox}
\vspace{1.1em}
Urutan \eng{fetch}--\eng{decode}--\eng{execute} tidak berarti BIU
dan EU selalu bekerja bergantian. Selama EU mengeksekusi instruksi
yang sudah diterjemahkan, BIU dapat mengisi \eng{queue} dengan
instruksi berikutnya. Karena itu, kedua unit tampak bekerja
bersamaan.
\end{infobox}

\begin{aktivitasbox}
\vspace{1.1em}
Perhatikan instruksi \lstasm|MOV AX, 1234h|.
\begin{enumerate}
    \item Tentukan unit yang mengambil instruksi tersebut dari
          memori.
    \item Tentukan unit yang menerjemahkan dan menjalankannya.
    \item Jelaskan mengapa hasil instruksi ini tidak memerlukan
          siklus bus tulis ke memori.
\end{enumerate}
\end{aktivitasbox}
